% Part: incompleteness
% Chapter: incompleteness-provability
% Section: lob-thm

\documentclass[../../../include/open-logic-section]{subfiles}

\begin{document}

\olfileid{inc}{inp}{tar}

\olsection{تعریف‌ناپذیری صدق}

مفهوم \emph{تعریف‌پذیری} به در اختیار داشتن معناشناسی صوری برای زبان
حساب وابسته است. مجموعه‌ای از فرمول‌ها و جمله‌ها را در زبان حساب توصیف
کرده‌ایم. «تفسیر مقصود» آن است که این جمله‌ها را بیانگر ادعاهایی دربارهٔ
اعداد طبیعی بخوانیم؛ چنین ادعایی می‌تواند صادق یا کاذب باشد. بگذارید
$\Struct{N}$ !!{structure}ی با دامنهٔ~$\Nat$ و تفسیر استاندارد
نمادهای زبان حساب باشد. آنگاه~$\Sat{N}{!A}$ یعنی «$!A$ در تفسیر
استاندارد صادق است».

\begin{defn}
رابطهٔ~$R(x_1,\dots,x_k)$ بر اعداد طبیعی در~$\Struct{N}$
\emph{تعریف‌پذیر} است اگر و تنها اگر فرمولی چون
$!A(x_1,\dots,x_k)$ در زبان حساب وجود داشته باشد به‌گونه‌ای که برای
هر~$n_1,\dots,n_k$، رابطهٔ~$R(n_1,\dots,n_k)$ برقرار باشد اگر و تنها
اگر~$\Sat{N}{!A(\num n_1,\dots,\num n_k)}$.
\end{defn}

به بیان دیگر، یک رابطه در~$\Struct{N}$ تعریف‌پذیر است اگر و تنها اگر
در نظریهٔ~$\Th{TA}$ بازنمایی‌پذیر باشد، که در آن
$\Th{TA} = \Setabs{!A}{\Sat{N}{!A}}$ مجموعهٔ جمله‌های صادق حساب است.
(اگر این مطلب بی‌درنگ برایتان روشن نیست، باید به تعریف‌ها بازگردید،
آن‌ها را بررسی کنید و خود را قانع سازید که چنین است.)

\begin{lem}
هر رابطهٔ محاسبه‌پذیر در~$\Struct{N}$ تعریف‌پذیر است.
\end{lem}

\begin{proof}
به‌آسانی می‌توان بررسی کرد فرمولی که رابطه‌ای را در~$\Th{Q}$ بازنمایی
می‌کند همان رابطه را در~$\Struct{N}$ تعریف می‌کند.
\end{proof}

اکنون می‌توان پرسید آیا عکس این نیز صادق است؟ یعنی آیا هر رابطه‌ای که
در~$\Struct{N}$ تعریف‌پذیر است محاسبه‌پذیر است؟ پاسخ منفی است. برای
نمونه:

\begin{lem}
رابطهٔ توقف در~$\Struct{N}$ تعریف‌پذیر است.
\end{lem}

\begin{proof}
به یاد آورید که قضیهٔ صورت نرمال کلینی می‌گوید هر تابع محاسبه‌پذیر
جزئی~$f$ نمایه‌ای چون~$e$ دارد به‌گونه‌ای که
$f(x) = U(\umin{s}{T(e,x,s)})$ برای هر~$x \in \Nat$، که در آن~$U$ و~$T$ بازگشتی اولیه و
بنابراین تام‌اند. پس~$f(x)$ تعریف شده است (یعنی محاسبه متوقف می‌شود)
اگر و تنها اگر~$s$ای وجود داشته باشد که~$T(e,x,s)$ برقرار باشد.

اکنون بگذارید~$H$ رابطهٔ توقف باشد، یعنی
\[
H = \Setabs{\tuple{e,x}}{\lexists[s][T(e, x, s)]}.
\]
بگذارید~$!D_T$، $T$ را در~$\Struct{N}$ تعریف کند. آنگاه
\[
H = \Setabs{\tuple{e,x}}{\Sat{N}{\lexists[s][!D_T(\num e, \num x, s)]}},
\]
پس~$\lexists[s][!D_T(z, x, s)]$ رابطهٔ~$H$ را در~$\Struct{N}$ تعریف
می‌کند.
\end{proof}

\begin{prob}
نشان دهید~$Q(n) \defiff n \in \Setabs{\Gn{!A}}{\Th{Q} \Proves !A}$
  در حساب تعریف‌پذیر است.
\end{prob}

خود~$\Th{TA}$ چطور؟ آیا در حساب تعریف‌پذیر است؟ یعنی آیا مجموعهٔ
$\Setabs{\Gn{!A}}{\Sat{N}{!A}}$ در حساب تعریف‌پذیر است؟ قضیهٔ تارسکی
پاسخ منفی می‌دهد.

\begin{thm}
\ollabel{thm:tarski}
مجموعهٔ !!{sentence}s صادق حساب در حساب تعریف‌پذیر نیست.
\end{thm}

\begin{proof}
فرض کنید~$!D(x)$ آن را تعریف کند؛ یعنی~$\Sat{N}{!A}$ اگر و تنها اگر
$\Sat{N}{!D(\gn{!A})}$. بنا بر لم نقطهٔ ثابت، جمله‌ای چون~$!A$
وجود دارد به‌گونه‌ای که
$\Th{Q} \Proves !A \liff \lnot !D(\gn{!A})$، و بنابراین
$\Sat{N}{!A \liff \lnot !D(\gn{!A})}$. اما آنگاه~$\Sat{N}{!A}$ اگر و
تنها اگر~$\Sat{N}{\lnot !D(\gn{!A})}$، و این با این واقعیت تناقض دارد
که~$!D(y)$ قرار است مجموعهٔ گزاره‌های صادق حساب را تعریف کند.
\end{proof}

تارسکی این تحلیل را بر مفهوم فلسفی عام‌تری از صدق اعمال کرد. تارسکی
استدلال کرد که برای هر زبان~$L$، مفهوم بسنده‌ای از صدق برای~$L$ باید
برای هر جملهٔ~$X$ شرط زیر را برآورده کند:
\begin{quote}
«$X$» صادق است اگر و تنها اگر~$X$.
\end{quote}
مثال بسیار نقل‌شدهٔ تارسکی برای زبان انگلیسی این جمله است:
\begin{quote}
«برف سفید است» صادق است اگر و تنها اگر برف سفید باشد.
\end{quote}
بااین‌حال، برای هر زبانی که به‌اندازهٔ کافی قوی باشد تا تابع قطری را
بازنمایی کند، و هر محمول زبانی~$T(x)$، می‌توانیم جمله‌ای چون~$X$ بسازیم
که شرط «$X$ اگر و تنها اگر نه~$T(\text{«$X$»})$» را برآورده کند. چون
نمی‌خواهیم محمول صدق برخی جمله‌ها را هم صادق و هم کاذب اعلام کند، تارسکی
نتیجه گرفت که نمی‌توان محمول صدقی برای همهٔ جمله‌های یک زبان مشخص کرد
مگر آنکه به‌نحوی از مرزهای آن زبان بیرون برویم. به بیان دیگر، محمول صدق
برای یک زبان را نمی‌توان در خود آن زبان تعریف کرد.

\end{document}
